% === Begin Label Data ===
% ID: 498
% 难度星级: 
% 标签: 
% 备注: 
% 组卷引用次数: 0
% === End  Label Data ===

\begin{problem}{2020}{G}{新高考I卷}{4}{三角函数}
日晷是中国古代用来测定时间的仪器，利用与晷面垂直的晷针投射到晷面的影子来测定时间．把地球看成一个球（球心记为 $ O $），地球上一点 $ A $ 的纬度是 $ OA $ 与地球赤道所在平面所成角，点 $ A $ 处的水平面是指过点 $ A $ 且与 $ OA $ 垂直的平面．在点 $ A $ 处放置一个日晷，若晷针与赤道所在平面平行，点 $ A $ 处的纬度为北纬 $ 40^\circ $，则晷针与点 $ A $ 处的水平面所成角为
\begin{choices}
\choice{{${20^\circ}$}}
\choice{{${40^\circ}$}}
\choice{{${50^\circ}$}}
\choice{{${90^\circ}$}}
\end{choices}
\end{problem}

\begin{answer}

\end{answer}

\begin{solutions}

\end{solutions}